\begin{problem}{Kaprekar Numbers}
    A natural number \(n\) is called a \textit{Kaprekar number} if there exists a way to split the decimal representation of \(n^2\) into a left part \(L\) and a right part \(R\)
    (where \(R\) may have leading zeros but must be positive) such that \(L + R = n\).
    For example:
    \begin{itemize}
      \item \(45^2 = 2025\), and \(20 + 25 = 45\),
      \item \(297^2 = 88209\), and \(88 + 209 = 297\),
      \item \(999^2 = 998001\), and \(998 + 001 = 999\).
    \end{itemize}
    Compute the sum of all Kaprekar numbers \(n\) under \(10^6\).

  
\end{problem}
